\documentclass[11pt,a4paper]{article}

\usepackage[utf8]{inputenc}
\usepackage[T1]{fontenc}
\usepackage{lmodern}
\usepackage[margin=2.2cm]{geometry}
\usepackage{amsmath}
\usepackage{booktabs}
\usepackage{array}
\usepackage{graphicx}
\usepackage{xcolor}
\usepackage{hyperref}
\usepackage{microtype}
\usepackage{caption}
\usepackage{subcaption}

\hypersetup{
  colorlinks=true,
  linkcolor=blue!55!black,
  citecolor=blue!55!black,
  urlcolor=blue!55!black,
  pdftitle={Chronological Benchmarking of Tabular Regression Models for Paris Apartment Valuation},
  pdfauthor={Badr Ettousy}
}
\graphicspath{{../../visuals/benchmark/}}
\setlength{\parskip}{0.35em}
\setlength{\parindent}{0pt}

\title{Chronological Benchmarking of Tabular Regression Models\\
for Paris Apartment Valuation Using DVF Transactions, 2021--2025}
\author{Badr Ettousy}
\date{August 2026}

\begin{document}
\maketitle

\begin{abstract}
This study benchmarks six methods for automated valuation of ordinary Paris
apartments using geolocated French Demandes de valeurs fonci\`eres (DVF)
transactions from 2021 through 2025. After mutation-level filtering and fixed
quality controls, the dataset contains 143,009 sales. Model selection uses
90,275 observations from 2021--2023 for development training and 24,404 sales
from 2024 for chronological validation. Selected configurations are refitted on
114,679 observations from 2021--2024 and evaluated on a common 2025 cohort of
28,330 transactions. The comparison includes an arrondissement median-price-per-
square-metre heuristic, ridge regression, random forest, histogram gradient
boosting, XGBoost, and LightGBM. LightGBM obtains the lowest 2025 mean absolute
error (MAE), EUR 96,981, with $R^2=0.8619$ and a median absolute percentage error
of 12.33\%. XGBoost is statistically indistinguishable at the 5\% level under a
paired bootstrap: the LightGBM-minus-XGBoost MAE difference is EUR $-387$ with a
95\% interval of approximately EUR $[-755,1]$. All four nonlinear tree ensembles
substantially outperform the market heuristic, whereas ridge regression performs
poorly. The results support gradient-boosted trees as the practical reference
class for this feature set, but do not establish that LightGBM universally
dominates XGBoost or future market regimes.
\end{abstract}

\textbf{Keywords:} automated valuation model, real estate, DVF, temporal
validation, gradient boosting, LightGBM, XGBoost, Paris.

\section{Introduction}

Automated valuation models (AVMs) estimate property transaction values from
structured characteristics and market evidence. They are attractive because
they are reproducible and scalable, yet they are vulnerable to spatial
heterogeneity, missing property attributes, heavy-tailed prices, and temporal
market shifts. A random train--test split can conceal the last problem because
near-contemporaneous transactions appear on both sides of the split.

This work asks a focused question: under a shared chronological protocol, which
of six common tabular-regression approaches provides the strongest out-of-year
valuation accuracy for Paris apartments? The contribution is empirical rather
than algorithmic. It consists of (i) mutation-aware DVF cleaning, (ii) a fixed
2021--2025 temporal evaluation design, (iii) model selection using a common
criterion, (iv) paired uncertainty estimates on identical 2025 transactions,
and (v) reproducible accuracy and computational-efficiency artifacts.

\section{Data}

\subsection{Source and scope}

DVF is produced by the French Direction g\'en\'erale des Finances publiques and
contains transaction information derived from notarial deeds and cadastral
records \cite{dvf2026}. This study uses the geolocated Paris files for calendar
years 2021--2025. The raw files contain 420,058 property rows representing
207,365 mutations.

DVF records a mutation-level value on one or more property rows. Treating every
row as an independent sale can therefore duplicate targets or associate a single
price with incompatible property types. The pipeline retains arm's-length sales
(``Vente'') containing exactly one apartment (DVF local-type code 2), permits
bundled dependencies such as a cellar or parking space (code 3), and rejects
mutations containing a house or commercial property (codes 1 or 4).

The remaining observations must contain a valid date, transaction value,
surface, Paris postcode, latitude, and longitude. Fixed plausibility bounds are
9--300 m$^2$ for surface, EUR 50,000--10,000,000 for value, and EUR
2,000--30,000 per m$^2$. Coordinates are restricted to the Paris bounding box
used by the application. Table~\ref{tab:data} reports the resulting cohorts.

\begin{table}[htbp]
\centering
\caption{Annual DVF row counts and eligible apartment transactions.}
\label{tab:data}
\begin{tabular}{lrrr}
\toprule
Year & Raw rows & Single-apartment sale rows & Clean model rows \\
\midrule
2021 & 81,517 & 31,835 & 30,397 \\
2022 & 98,308 & 34,566 & 33,080 \\
2023 & 80,993 & 28,174 & 26,798 \\
2024 & 74,800 & 25,595 & 24,404 \\
2025 & 84,440 & 29,503 & 28,330 \\
\midrule
Total & 420,058 & 149,673 & 143,009 \\
\bottomrule
\end{tabular}
\end{table}

\subsection{Features and target}

The target is total mutation value in euros. Learned models optimize
$\log(1+y)$ and transform predictions back with $\exp(\hat{z})-1$, reducing the
influence of the long upper tail while preserving a positive-scale prediction.
The feature vector contains built surface, number of principal rooms, latitude,
longitude, address number, number of lots, month, day of year, calendar year,
months since January 2021, and postcode. Numeric missing values are median
imputed; postcode is one-hot encoded. Numeric inputs are standardized only for
ridge regression, where scale directly affects the penalty.

Important determinants are unavailable in DVF: interior condition, floor,
elevator, energy rating, exposure, view, occupancy, and detailed building
quality. The estimated error therefore reflects both model limitations and
irreducible missing-information error.

\section{Methods}

\subsection{Temporal evaluation protocol}

The split is chronological and identical for every method:
\begin{itemize}
  \item development training: 2021--2023 (90,275 rows);
  \item configuration selection: 2024 (24,404 rows), minimizing MAE;
  \item final refit: 2021--2024 (114,679 rows);
  \item common test: 2025 (28,330 rows).
\end{itemize}

No 2025 transaction is used by a fitting or configuration-selection call in the
benchmark program. However, 2025 results from an earlier baseline experiment had
already been inspected by the analyst. Consequently, 2025 is computationally
held out for the six-model comparison but should not be described as a fully
blind external test.

\subsection{Models}

The median-price baseline multiplies apartment surface by the median price per
m$^2$ observed in its arrondissement during the available training years. Ridge
regression adds an $\ell_2$ penalty to a linear model \cite{hoerl1970}. Random
forest averages decorrelated trees trained on bootstrap samples and randomized
feature subsets \cite{breiman2001}. Histogram Gradient Boosting is the existing
project model and represents conventional histogram-binned boosting
\cite{friedman2001}. XGBoost adds a regularized, systems-optimized tree-boosting
implementation \cite{chen2016}; LightGBM uses histogram learning and leaf-wise
growth with efficiency techniques including gradient-based one-side sampling and
exclusive feature bundling \cite{ke2017}.

Table~\ref{tab:config} gives the selected configuration. Each learned model is
chosen from a compact, predefined search space of three candidates, except ridge
regression, which uses four penalty values. This is a controlled engineering
benchmark, not an exhaustive hyperparameter optimization.

\begin{table}[htbp]
\centering
\caption{Configurations selected by 2024 validation MAE. Shared random seed is 42.}
\label{tab:config}
\small
\begin{tabular}{p{3.4cm}p{10.2cm}}
\toprule
Method & Selected configuration \\
\midrule
Median price/m$^2$ & Arrondissement median from the training period \\
Ridge & $\alpha=100$ \\
Random Forest & 240 trees; max-features 0.7; minimum leaf size 5 \\
Histogram Gradient Boosting & 63 leaves; minimum leaf size 40; $\ell_2=2$; learning rate 0.06 \\
XGBoost & depth 6; minimum child weight 5; 650 trees; learning rate 0.04 \\
LightGBM & 31 leaves; minimum child size 25; 650 trees; learning rate 0.04 \\
\bottomrule
\end{tabular}
\end{table}

\subsection{Evaluation criteria}

For actual values $y_i$ and predictions $\hat{y}_i$, the primary criterion is
\begin{equation}
\mathrm{MAE}=\frac{1}{n}\sum_{i=1}^{n}|y_i-\hat{y}_i|.
\end{equation}
Secondary criteria are root mean squared error (RMSE), coefficient of
determination ($R^2$), mean absolute percentage error (MAPE), median absolute
percentage error (MdAPE), and the proportions within 10\% and 20\% of the actual
value. Training time, batch prediction latency, and serialized artifact size are
also recorded.

Individual MAE intervals use 500 nonparametric bootstrap resamples. Model-to-
model conclusions use 1,000 paired bootstrap resamples of per-transaction
absolute-error differences, preserving the fact that all models predict the same
sales. A paired interval wholly below zero indicates that the row model has lower
MAE at the 5\% level under this resampling procedure.

\section{Results}

\subsection{Overall accuracy}

Table~\ref{tab:results} presents the common 2025 test results. LightGBM has the
lowest observed MAE at EUR 96,981, followed by XGBoost at EUR 97,367 and
Histogram Gradient Boosting at EUR 97,753. The same three models are almost tied
on typical relative error, with MdAPE between 12.33\% and 12.48\%. LightGBM puts
71.41\% of test predictions within 20\% of the transaction value, compared with
62.60\% for the market baseline.

\begin{table}[htbp]
\centering
\caption{Performance on 28,330 transactions from the common 2025 test cohort.}
\label{tab:results}
\resizebox{\textwidth}{!}{%
\begin{tabular}{lrrrrrrrr}
\toprule
Model & MAE (EUR) & 95\% MAE interval & RMSE (EUR) & $R^2$ & MdAPE & Within 10\% & Within 20\% \\
\midrule
LightGBM & 96,981 & [94,935; 99,014] & 197,107 & 0.8619 & 12.33\% & 41.49\% & 71.41\% \\
XGBoost & 97,367 & [95,312; 99,485] & 199,023 & 0.8592 & 12.34\% & 41.22\% & 71.32\% \\
Hist. Gradient Boosting & 97,753 & [95,730; 99,884] & 198,027 & 0.8606 & 12.48\% & 40.87\% & 70.87\% \\
Random Forest & 99,797 & [97,726; 101,900] & 200,909 & 0.8566 & 12.89\% & 40.20\% & 70.06\% \\
Median price/m$^2$ & 110,868 & [108,575; 113,272] & 221,015 & 0.8264 & 14.62\% & 36.54\% & 62.60\% \\
Ridge regression & 169,126 & [163,427; 174,924] & 549,280 & $-0.0722$ & 19.99\% & 25.51\% & 50.03\% \\
\bottomrule
\end{tabular}}
\end{table}

\begin{figure}[htbp]
\centering
\includegraphics[width=\textwidth]{benchmark_accuracy.png}
\caption{Accuracy comparison, operational hit rate, and validation-to-test transfer. Error bars in the MAE panel are 95\% bootstrap intervals.}
\label{fig:accuracy}
\end{figure}

Relative to the median-price baseline, LightGBM reduces test MAE by EUR 13,887
(paired 95\% interval EUR 12,896--14,905), a 12.53\% reduction. The corresponding
reductions are EUR 13,501 for XGBoost, EUR 13,115 for Histogram Gradient
Boosting, and EUR 11,071 for Random Forest; their paired intervals against the
baseline are all strictly positive.

\subsection{Are the leading models distinguishable?}

The LightGBM-minus-XGBoost paired MAE difference is EUR $-387$ with a 95\%
interval of EUR $[-755,1]$. Thus, the observed ordering does not support a
5\%-level claim that LightGBM is more accurate than XGBoost. LightGBM is more
accurate than Histogram Gradient Boosting by EUR 772 per sale, with interval EUR
$[400,1{,}140]$ in LightGBM's favor, and more accurate than Random Forest by EUR
2,816, with interval EUR $[2{,}257,3{,}331]$. Figure~\ref{fig:pairwise} reports
the full paired difference matrix.

\begin{figure}[htbp]
\centering
\includegraphics[width=0.82\textwidth]{benchmark_pairwise_differences.png}
\caption{Mean paired absolute-error differences in thousands of euros. Negative values favor the row model.}
\label{fig:pairwise}
\end{figure}

\subsection{Efficiency and temporal stability}

On the benchmark machine, the selected LightGBM refit took 2.05 seconds and its
serialized artifact occupied 1.84 MB. XGBoost took 3.27 seconds and occupied
2.60 MB; Histogram Gradient Boosting took 2.87 seconds and occupied 1.36 MB.
Random Forest required 16.63 seconds and 383.27 MB. These single-run wall-clock
measurements describe this implementation and hardware, not universal library
speed rankings.

\begin{figure}[htbp]
\centering
\includegraphics[width=\textwidth]{benchmark_efficiency.png}
\caption{Observed accuracy--efficiency trade-offs. Training-time axes are logarithmic; bubble area represents serialized artifact size.}
\label{fig:efficiency}
\end{figure}

All nonlinear ensembles retain a clear advantage over the baseline throughout
2025. The boosted models have monthly MdAPE near 12\% in the first half and
generally rise toward 13--14\% by December, suggesting modest late-year drift.
Ridge is consistently weakest and deteriorates more strongly in the second half.

\begin{figure}[htbp]
\centering
\includegraphics[width=\textwidth]{benchmark_monthly_stability.png}
\caption{Month-by-month median absolute percentage error on the 2025 test cohort.}
\label{fig:monthly}
\end{figure}

\section{Discussion}

The most important result is class-level rather than library-level: nonlinear
tree ensembles are clearly preferable to the market heuristic and the global
linear model for the available DVF features. Location coordinates, postcode,
surface, and their interactions create irregular valuation functions that ridge
regression cannot represent without explicit basis expansion. Its negative test
$R^2$ and high RMSE also indicate sensitivity to expensive-tail transactions
after retransformation from log space.

Among tree ensembles, gradient boosting is the strongest family. LightGBM is the
best point estimate and offers an attractive accuracy--size--training-time
combination. Nevertheless, its difference from XGBoost is too small to declare a
robust accuracy victory on this test cohort. A production choice between the two
should therefore consider operational simplicity, explainability tooling,
latency under the actual service workload, and monitoring behavior in newer
market periods.

The strong arrondissement median-price baseline is also informative. It reaches
$R^2=0.8264$ using only surface and arrondissement-level market evidence. The
boosted models improve MAE by approximately 12\%--13\%, demonstrating genuine
incremental value, but the remaining EUR 97k MAE shows that the current feature
set is not sufficient for appraisal-grade precision.

\section{Limitations and validity threats}

First, the 2025 cohort is temporally held out by the benchmark code but was
previously inspected during development of the original project model. It is not
an untouched external replication. A future publication claim should be tested
on a newly released period or another geography that has not informed model or
feature decisions.

Second, DVF does not contain many causal value drivers, and mutation-level prices
may include dependencies. The eligibility rules reduce ambiguity but do not
recover unobserved condition or amenity information. Third, the search is compact
and slightly different in size across model families. It is designed to avoid a
large, unfair tuning advantage, but it does not locate each method's global
optimum. Fourth, bootstrap intervals condition on the observed 2025 transaction
population and do not quantify regime-change uncertainty. Finally, timing values
come from one Windows workstation and one run per selected final model.

\section{Reproducibility}

The experiment uses Python 3.12, NumPy 1.26.4, pandas 2.2.2, scikit-learn 1.5.1,
XGBoost 3.4.0, LightGBM 4.7.0, and joblib 1.4.2. Random seeds are fixed at 42 for
model fitting and 20250811 for bootstrap resampling. The complete benchmark took
294.1 seconds in the recorded run. The project stores the cleaning and benchmark
code, annual input manifest, selected model artifacts, per-row validation and
test predictions, hyperparameter-search results, pairwise comparisons, and all
figures required by this paper.

The principal commands are:
\begin{verbatim}
python benchmark_models.py
python visualize_benchmark.py
powershell -ExecutionPolicy Bypass -File docs/paper/compile_paper.ps1
\end{verbatim}

\section{Conclusion}

Under a shared chronological evaluation, LightGBM produces the lowest observed
2025 MAE, while XGBoost is statistically tied under paired resampling. Histogram
Gradient Boosting and Random Forest remain competitive, and all four nonlinear
ensembles outperform the arrondissement median-price baseline. Ridge regression
is unsuitable as the principal AVM with the current representation. The next
scientific step is not to assume that a more complex architecture will win, but
to add richer building and neighborhood information and benchmark any attention-
based model against the LightGBM/XGBoost reference under a newly locked temporal
or geographic test.

\section*{Acknowledgements}
The computational workflow and manuscript preparation were assisted by an AI
coding system. Responsibility for data use, analysis choices, verification, and
all scientific claims remains with the project author.

\bibliographystyle{plain}
\bibliography{references}

\end{document}
